% ─────────────────────────────────────────────────────────────────────────────
\section{Training Dataset Construction}
\label{app:dataset}
% ─────────────────────────────────────────────────────────────────────────────

This appendix details the construction of the 82K image--prompt corpus used
to train the Auditor-Scorer $\As$, covering source datasets, prompt
reconstruction, labeling pipeline design, and labeling quality.

% ─────────────────────────────────────────────────────────────────────────────
\subsection{Source Datasets}
\label{app:dataset_sources}
% ─────────────────────────────────────────────────────────────────────────────

We draw from eight publicly available HuggingFace datasets spanning three
content categories: safe, nudity, and violence.
Table~\ref{tab:dataset} summarizes each source, its nominal category,
and whether prompts were available or reconstructed.

\begin{table*}[!ht]
\centering
\small
\setlength{\tabcolsep}{5pt}
\renewcommand{\arraystretch}{1.15}
\caption{Source datasets for the 82K auditor training corpus from huggingface.}
\label{tab:dataset}
\begin{tabular}{lcc}
\toprule
\textbf{Dataset} & \textbf{Category} & \textbf{Prompts} \\
\midrule
\texttt{Subh775/WeaponDetection} \cite{yolov7test_weapon_detection}              & Violence         & VLM-generated \\
\texttt{NeuralShell/Gore-Blood-Dataset-v1.0} \cite{Gore-Blood-Dataset-v1.0}  & Violence         & VLM-generated \\
\texttt{x1101/nsfw-full}                      & Nudity           & VLM-generated \\
\texttt{Lenkashell/unsafe\_violence\_image\_captions}  & Violence  & Existing \\
\texttt{Lenkashell/unsafe\_shocking\_image\_captions}  & Violence  & Existing \\
\texttt{yiting/UnsafeBench}\cite{10.1145/3719027.3765088}                   & Nudity, Violence & Existing \\
\texttt{OpenSafetyLab/t2i\_safety\_dataset}\cite{Li_2025_CVPR}   & Nudity, Violence & Existing \\
\bottomrule
\end{tabular}
\end{table*}

Sources were selected to cover both the nudity and violence policy axes with
sufficient volume and diversity of visual style, scene complexity, and
severity level. Datasets with pre-existing captions were used as-is after
cleaning; datasets without prompts required reconstruction (see
Section~\ref{app:prompt_reconstruction}).

% ─────────────────────────────────────────────────────────────────────────────
\subsection{Prompt Reconstruction}
\label{app:prompt_reconstruction}
% ─────────────────────────────────────────────────────────────────────────────

Four source datasets did not include generation prompts. For these, image
captions were synthesized using Qwen2.5-VL~\cite{bai2025qwen25vltechnicalreport}
as a surrogate prompt. The model was instructed to produce a concise,
descriptive caption in the style of a text-to-image generation prompt rather
than a natural-language description, targeting the visual content of the image
without editorializing about policy compliance. These reconstructed captions
serve the faithfulness branch of the auditor during training and are not used
for any safety label decision.

% ─────────────────────────────────────────────────────────────────────────────
\subsection{Labeling Pipeline}
\label{app:labeler_ablation}
% ─────────────────────────────────────────────────────────────────────────────

All images---regardless of source label---were relabeled into three mutually
exclusive categories: $\{\mathrm{safe},\,\mathrm{nudity},\,\mathrm{violence}\}$.
Multi-label annotation was evaluated and discarded: overlapping soft gradients
from multi-label supervision degraded auditor classification performance in
preliminary experiments. One-hot categorization was therefore adopted as the
final labeling scheme.

\paragraph{Labeler selection.}
We evaluated five candidate VLMs as automated labelers. Each was assessed on
a manually inspected reference set of hard cases including bikinis, underwear,
sports scenes with physical contact, crowds, close interpersonal proximity,
and classical artwork (e.g., Greek nude sculpture). The reference set was
chosen to stress-test policy-boundary behavior rather than gross violations,
which all models handle trivially.

Table~\ref{tab:labeler_ablation} summarizes observed failure modes.

\begin{table}[h]
\centering
\small
\setlength{\tabcolsep}{4pt}
\renewcommand{\arraystretch}{1.2}
\caption{Labeler ablation: observed failure modes on the hard-case calibration
set. FP = false positive rate; FN = false negative rate on the calibration
reference.}
\label{tab:labeler_ablation}
\begin{tabular}{p{0.22\linewidth}p{0.68\linewidth}}
\toprule
\textbf{Model} & \textbf{Observed failure mode} \\
\midrule
Qwen2.5-VL (7B) &
  Basic prompting yields high FP rate; chain-of-thought prompting introduces
  label contradictions; self-consistency is slow and inconsistent;
  few-shot prompting degrades prompt-following according to community
  reports. \\
\addlinespace
Qwen3-VL &
  Similar failure pattern to Qwen2.5-VL; no improvement on boundary cases. \\
\addlinespace
ShieldGemma (vanilla) &
  Slow inference; inconsistent categorization on borderline cases without
  fine-tuning; high FN rate overall. \\
\addlinespace
Grounding DINO~\cite{liu2023gdino} &
  Highly threshold-sensitive (small shift produces either high FP or high
  FN); lacks semantic context for borderline cases (\eg, classical nude
  sculpture). \\
\addlinespace
\textbf{InternVL-3.5 (8B)} &
  \textbf{Best overall performance}: lowest FP rate; strong contextual
  nuance (\eg, correctly categorizes classical nude sculptures as safe);
  conservative on FN (acceptable for training corpus construction). \\
\bottomrule
\end{tabular}
\end{table}

InternVL-3.5 (8B) was selected as the final labeler based on this ablation.

\paragraph{Why chain-of-thought prompting was de-emphasized.}
Explicit chain-of-thought prompting increases verbosity without reliably
improving category fidelity. On ambiguous images, it causes models to
generate unsupported post-hoc rationales or self-contradictory
intermediate reasoning that conflicts with the final label.
Self-consistency (repeated passes with majority voting) was also rejected: it
increases inference cost substantially while failing to correct systematic
classification errors caused by stable policy-boundary confusion---the primary
failure mode on borderline cases.

\paragraph{Policy-first prompt design.}
The InternVL-3.5 prompt was designed to be policy-oriented rather than
explanation-oriented. Categories are defined in operational terms directly tied
to visual evidence:

\begin{itemize}[leftmargin=1.2em,itemsep=2pt,topsep=2pt]
  \item \textbf{Nudity:} Swimwear, underwear,
        medical imagery, and athletic clothing are explicitly safe.
  \item \textbf{Violence:} visible physical harm or graphic cues such as
        blood or weapons in active use. Proximity, sparring poses, and
        non-graphic contact are safe.
  \item \textbf{Safe:} all other content, including classical artwork with
        nudity, sports imagery without graphic harm, and intimate proximity
        without explicit content.
\end{itemize}

This framing prevents surface-level cue triggering (e.g., visible skin area,
figure proximity, athletic posture) which was the dominant failure pattern
across rejected labelers.

% ─────────────────────────────────────────────────────────────────────────────
\subsection{Threshold Calibration and Dead-Zone Exclusion}
\label{app:threshold_calibration}
% ─────────────────────────────────────────────────────────────────────────────

InternVL-3.5 produces a continuous policy score for each image. Because no
large ground-truth labeled set was available for data-driven threshold
selection, the operating threshold was calibrated empirically against the
same manually inspected hard-case reference set used for labeler ablation.
Model scores were compared against human judgment on this subset, and the
threshold was set to reflect the desired moderation policy rather than an
externally optimized operating point.

A \textbf{dead zone} was reserved around the threshold: images whose scores
fell within a margin of the decision boundary were excluded from the training
corpus entirely rather than assigned an arbitrary label. This exclusion
strategy ensures that the high-confidence tails of the score distribution---
where label quality is highest---dominate training, while uncertain borderline
samples do not introduce spurious supervision signal.

\paragraph{Residual label noise.}
The most likely residual failure mode from single-shot VLM scoring is not
catastrophic random mislabeling but systematic boundary bias: InternVL-3.5
may mildly over-score skin-heavy but policy-safe images (e.g., bikini
photographs), causing a fraction of such cases to drift above the nudity
threshold. This produces moderate boundary noise in the minority class rather
than label collapse. For auditor training, such boundary noise is manageable
provided the high-confidence tails remain clean---which the dead-zone
exclusion strategy is designed to ensure.

% ─────────────────────────────────────────────────────────────────────────────
\subsection{Dataset Statistics and Splits}
\label{app:dataset_stats}
% ─────────────────────────────────────────────────────────────────────────────

The final labeled corpus contains 82,000 image--prompt pairs distributed
across the three categories. After dead-zone exclusion of boundary cases, the
corpus was split into training, validation, and test sets using a
70\,/\,15\,/\,15 stratified partition, preserving class proportions across
splits. The test set was held out entirely for the classification evaluation
reported in Table~\ref{tab:auditor_perf}.